% !TeX spellcheck = en_US
\chapter{Discussion}\label{ch:discussion}%

This chapter interprets the results of Chapters~\ref{ch:sim} and~\ref{ch:real} with respect to the four questions of Section~\ref{sec:research-questions}, relates the offline and online distillation variants, discusses the action and observation contracts, returns to the questions of the exposé, and states the limitations of the work.

\section{Interpretation of the Ablation Study}\label{sec:disc-ablations}

\paragraph{Unified multi-cycle student (Q1).}
The single-cycle student and the multi-cycle student share the same architecture, the same encoder, the same real-world augmentation, and the same per-cycle behavior; they differ only in whether the training data spans several cycles and in whether the temporal context is reset between cycles. The multi-cycle student completes every trial in both domains, the single-cycle student about half of the cycles. The failure mode, slipping off the nut edge after a few cycles, identifies what the single-cycle data lacks: the axial advance of the nut. In a single-cycle demonstration the nut height at grasp is always the reset height, so the student never has to read the height from the point cloud and adjust its approach; in a multi-cycle demonstration the height changes by one sixth of a pitch per cycle and the student must track it. The global progress scalar supports this, but it is not the cause of the difference, because the single-cycle student also receives it; the cause is the coverage of the axial state in the demonstrations. The result supports the central design principle: it is the unification of the complete cyclic behavior in one dataset, not a per-cycle controller invoked repeatedly, that yields a policy that sustains consecutive cycles.

\paragraph{Point-cloud encoder (Q2).}
The pose probe shows where the image-space representation falls short: its positional accuracy equals that of the point-cloud representation, but its rotational error is twice as large. Re-grasping requires the orientation of the nut faces relative to the fingers; a representation that encodes position but not orientation lets the fingers close on the right spot with the wrong face alignment, which is the error that the closed-loop ablation exhibits as a large positive cyclic error and early failure. The point-cloud encoder has no pretrained backbone and about 60\,000 parameters, yet it generalizes from the randomized simulated views to the physical camera, whereas the ResNet student degrades further on hardware. Two reasons are plausible and not separated by the experiments: the point encoder consumes metric geometry that is identical in both domains up to depth noise, and random point sampling acts as a strong augmentation that the image encoder does not receive. The pretraining experiments add a caveat: initializing the encoder from a pose-regression pretraining and then training it jointly with the flow objective moved the latent so that the pretrained decoder no longer applied, and the resulting student was worse than one trained from the sweep initialization alone. The visual feature that the flow objective shapes is not a pose estimate, and attempts to make it one through pretraining should either freeze the encoder or supervise orientation directly.

\paragraph{Flow matching versus deterministic regression (Q3).}
Both regression heads fail in simulation, the recurrent one immediately and the chunked Transformer after about one cycle, and both fail by compounding error. The chunked Transformer regression head shares the backbone, the chunk length, and the temporal ensembling of the flow student, so the difference between (d) and (f) isolates the generative objective. Two mechanisms explain the gap. The first is multimodality of the labels: the dataset contains, for near-identical observations, the teacher's nominal action, the teacher's recovery after an impulse, and the operator's correction; a squared-error head averages them and produces an action that none of the sources would take, whereas the flow head samples one of them. The second is the structure of the prediction: the flow head predicts a velocity toward the data manifold from a noisy chunk, which is a denoising problem that is easier to fit than a direct regression from observation to chunk and which, through the sixteen integration steps, refines the chunk jointly over all eight steps with full attention among them. The result does not imply that regression heads cannot thread; it implies that, at this data scale and with this heterogeneous data, the generative head is markedly more robust.

\paragraph{Human-corrected real-world data (Q4).}
The student trained on simulated data only is the best student in simulation and fails on hardware, completing two cycles with a cyclic error five times that of the full method. The 6.5\,\% of real frames that the operator overwrote, together with the 93.5\,\% of real frames that replay the teacher, are enough to close the gap. The mechanism is that the corrected segments cover exactly the states in which the simplified thread model and the real contact disagree, the grasp-to-release interval, with labels that succeed on hardware, and the uncorrected real frames supply the real camera and proprioception statistics for the rest of the cycle. The approach is a form of DAgger in which the real robot is the environment and the operator is the expert, restricted to a single round and to the segments the operator chooses. It requires no real-world reward and no real-world policy optimization, and its cost, twenty replays with sparse interventions, is small compared with the cost of real-world reinforcement learning.

\section{Offline versus Online Distillation}\label{sec:disc-online}

The thesis set out, in the exposé and in the code base, to distill the teacher online with DAgger, and it ends with an offline flow student. The comparison in Section~\ref{sec:sim-online} explains why. The online recurrent students learn either the turn or the transition depending on who controls the environment, and a decayed intervention schedule lets them learn both but not sustain them beyond two cycles. The evaluation ladder shows that these students fit their labels and execute them correctly under teacher-forced observations, so their failure is one of recovery under their own observations. The offline flow student recovers better because its training data contain explicit recoveries from disturbances, because its chunked generative prediction with temporal ensembling smooths single-step errors, and because the progress and phase conditions keep it synchronized with the physical cycle. DAgger's advantage, labels on student-visited states, is obtained here in a weaker form by the impulses, which perturb the teacher rather than the student; the online flow variant with $\beta=1$ obtains labels on teacher-visited states only and therefore reduces to online collection of the same data, and the chunked-label constraint prevents the mixture schedules that help the one-step recurrent student. The online variant with student-visited states would need multi-step teacher recovery fragments collected from student states, which the code base plans but does not yet implement. The practical argument is decisive for deployment: the real-world data is offline by nature, so an offline student is required in any case, and a student that trains from the mixed dataset without changing its objective is preferable to one that needs an online loop in simulation and a separate fine-tuning stage for real data.

\section{Action and Observation Contracts}\label{sec:disc-contracts}

The project implemented two action contracts and settled on the Cartesian one. The joint-PD contract was motivated by deployment convenience and by the expectation that joint references are easier for a controller to track than Cartesian targets, which they are: recorded joint references replayed at the physics rate reproduce the teacher motion with sub-millimeter accuracy. The contract fails as a learning target, not as a control interface. Converting the teacher's bounded Cartesian action into seventeen physical-unit references at the physics rate and then asking the student to emit one such reference per 15\,Hz step introduces a hold error that the teacher never produced, mixes units with scale differences of three orders of magnitude in one regression target, and removes the bolt-relative structure of the action that lets the student generalize across nut heights. Every offline student trained under this contract reproduced its labels and diverged in closed loop, in the full-task case within half a second. The Cartesian contract keeps the student in the teacher's own action space and executes it through the unchanged controller, including the action smoothing that the teacher was trained with. The thesis does not conclude that joint-space students are infeasible; it concludes that a student should imitate the action the teacher produced, in the units the teacher produced it, and let the controller do the conversion it was designed for.

The observation contract follows the same logic. The student receives the previous filtered action and the controlled joints, which is what the teacher's history contains, and the two low-dimensional conditions that the demonstration provides without privileged state. The global progress scalar is the weakest part of this contract: it assumes a nominally timed trial, and the long-horizon stress trials show that the student's internal timing drifts behind the teacher by several seconds after ten cycles. The process-phase encoding corrects this within a cycle, but a trial that runs much longer than the demonstrations clamps the progress at one and leaves the student without a clock. Inferring the phase from the observation history inside the policy, rather than from an external signal, is the natural next step and is discussed in Chapter~\ref{ch:conclusion}.

\section{Answers to the Questions of the Exposé}\label{sec:disc-expose}

The first question of the exposé asked whether vanilla PPO can learn a force-aware threading policy with a dexterous hand--arm system in a privileged, high-fidelity simulation without constraining the thread interaction. The answer from this work is negative in the form asked: without a constraint on the thread interaction, PPO exploited the SDF contact instead of threading, and with grasp acquisition, grasp retention, helical progress, and termination all shaped in one reward, no reward variant produced a threading policy. The answer is positive once the thread is constrained kinematically and the task is reduced to one turn from an engaged nut: PPO then learns a robust turn under initial-pose and controller randomization, and the hybrid teacher extends it to arbitrarily many cycles. Force awareness, the second part of the question, was not needed for the turn: the actor of the retained teacher observes no force channel, and hardware-equivalent force observations were implemented only at the end of the project and remain untested.

The second question asked whether the privileged teacher can be distilled into a vision-based student with the localization accuracy that sub-millimeter thread interaction requires. The answer is that the student does not need to localize the nut to sub-millimeter accuracy, because the kinematic engagement of the nut on the bolt and the compliance of the impedance controller absorb lateral error; what the student needs is the orientation of the faces for the re-grasp and the height of the nut for the approach, and the point-cloud representation provides both with a precision that completes ten cycles on hardware across three sizes. The exposé's expectation that explicit depth would be required is confirmed, and its expectation that resolution or stereo would be the deciding factor is not: the deciding factor is the geometric structure of the representation.

\section{Limitations}\label{sec:disc-limitations}

The work has limitations that the results make explicit. The nut starts engaged on the bolt; autonomous engagement from a loose nut was not achieved by the pickup expert, so the full task from table to tightened nut is not demonstrated. The demonstrated behavior is a pinch-grasp re-grasp strategy tied to the face topology of a hexagon; it transfers to profiles that preserve compatible re-grasp configurations and fails on the rectangle, and the student has no mechanism to discover a different grasp. The policy fails after 15 to 22 cycles because the axial range of its demonstrations is bounded, and the global progress clock does not extend beyond the nominal trial length. The evaluation uses ten trials per condition and reports means and standard deviations without confidence intervals, and several design ablations in Chapter~\ref{ch:sim} are matched single rollouts at one seed, which establish the direction of an effect but not its variance. The teacher's thread model is a kinematic simplification, and the sim-to-real gap it creates is closed by human corrections whose collection interface and selection criteria are the operator's judgment, which is not reproducible in the strict sense. Finally, the hardware experiments were run with a measured but not deployment-validated camera calibration, and the dependence of the student on calibration error has not been quantified.
